\documentclass[11pt]{article}

\usepackage[margin=1in]{geometry}
\usepackage{amsmath,amssymb,amsthm,mathtools,bm}
\usepackage{mathrsfs}
\usepackage{bbm}
\usepackage{microtype}
\usepackage{xcolor}

\title{Global Convergence for Low-Rank Mean-Field Neural Networks:\\Complete Proof for Arbitrary Depth}
\author{ }
\date{}

% math macros
\newcommand{\E}{\mathbb{E}}
\newcommand{\Pp}{\mathbb{P}}
\newcommand{\R}{\mathbb{R}}
\newcommand{\1}{\mathbbm{1}}
\newcommand{\ip}[2]{\left\langle #1,\,#2 \right\rangle}
\newcommand{\norm}[1]{\left\lVert #1 \right\rVert}
% double brackets for norms
\newcommand{\llbracket}{[\![}
\newcommand{\rrbracket}{]\!]}
% interleave for Orlicz norms
\newcommand{\interleave}[1]{\llbracket #1 \rrbracket}

\newtheorem{theorem}{Theorem}
\newtheorem{lemma}{Lemma}
\newtheorem{definition}{Definition}
\newtheorem{assumption}{Assumption}
\newtheorem{remark}{Remark}
\newtheorem{proposition}{Proposition}

\begin{document}

\maketitle

\tableofcontents

\section{Introduction}

This document provides a complete proof of global convergence for low-rank mean-field neural networks of arbitrary depth $L \ge 2$. The key advantage of our low-rank framework is that it achieves global convergence for arbitrary depth with standard i.i.d. initialization, unlike full-width networks which require correlated initialization for $L \ge 3$.

\subsection{Key Differences from Full-Width Case}

\begin{itemize}
\item \textbf{Frozen random features}: The first layer uses frozen random features $f_1(C_1)$, eliminating the need for algebraic topology arguments to maintain full support.
\item \textbf{No correlated initialization needed}: Standard i.i.d. initialization suffices for arbitrary depth $L$.
\item \textbf{Simpler proof structure}: The frozen first layer provides a stable base, making the proof structure cleaner than the full-width case.
\item \textbf{Polynomial scaling}: The Gronwall constant grows as $(1+rK)^{L-2}$ (polynomial in $L$), not exponentially.
\end{itemize}

\section{Setup: Low-Rank Mean-Field Equations}

We consider an $L$-layer low-rank random feature network with the following structure:

\begin{itemize}
\item \textbf{Layer 1}: Random features $\varphi_1(\ip{f_1(C_1)}{X})$ where $f_1(C_1)$ are \textbf{frozen} random vectors and $C_1$ indexes the first-layer channels. \textbf{This layer is not trained.}
\item \textbf{Layers $i = 2, \ldots, L-1$}: Each layer employs a low-rank mixing matrix $L^{(i)}$ of rank $r$, producing pre-activations
\[
H_i(t,c_i;X,W)=\sum_{k=1}^r L^{(i)}_{c_i,k}\,m_k^{(i)}(t;X,W),
\]
where $m_k^{(i)}(t;X,W)$ are the $r$ channel-wise partial functions at layer $i$.
\item \textbf{Layer $L$ (output)}: Trainable weights $w_L(t,C_L)$ producing the final output.
\end{itemize}

The mean-field ODE system consists of $(L-2)\times r + 1$ trainable equations:
\begin{align}
\partial_t w_1(t,c_1,k) &= -\xi_1(t)\,\E_{Z}\!\left[d_L\big(Z;W(t)\big)\,\varphi_1\!\left(\ip{f_1(c_1)}{X}\right)\,B_k^{(2)}(t;X)\right], \label{eq:ode-w1}\\
\partial_t w_i(t,c_i,k) &= -\xi_i(t)\,\E_{Z}\!\left[d_L\big(Z;W(t)\big)\,\varphi_i\!\Big(H_i(t,c_i;X,W(t))\Big)\,B_k^{(i+1)}(t;X)\right], \quad i=2,\ldots,L-1, \label{eq:ode-wi}\\
\partial_t w_L(t,c_L) &= -\xi_L(t)\,\E_{Z}\!\left[d_L\big(Z;W(t)\big)\,\varphi_L\!\Big(H_L(t,c_L;X,W(t))\Big)\right], \label{eq:ode-wL}
\end{align}
where $d_L(Z;W)=\partial_2 \mathcal{L}(Y,\hat y(X;W))$ is the loss derivative, and for $i=2,\ldots,L-1$:
\[
B_k^{(i)}(t;X)=\E_{C_i}\!\left[L^{(i)}_{C_i,k}\,\varphi_i'\!\Big(H_i(t,C_i;X,W(t))\Big)\,w_i(t,C_i,k)\right]
\]
is the channel-$k$ backpropagated signal at layer $i$.

The model output is:
\[
\hat y(X;W(t)) = \E_{C_L}\big[w_L(t,C_L)\,\varphi_L\big(H_L(t,C_L;X,W(t))\big)\big].
\]

\section{Assumptions for Global Convergence}

\begin{assumption}[Initialization]\label{assump:init-lowrank}
For each layer $\ell\in\{2,\dots,L\}$, the initial functions satisfy:
\begin{align*}
\sup_{m\ge 1}\frac{1}{\sqrt{m}}\max_{1\le k\le r}\E_{C_\ell}\!\left[|w_\ell^0(C_\ell,k)|^m\right]^{1/m} &\le K \quad \text{for } \ell\in\{2,\dots,L-1\},\\
\sup_{m\ge 1}\frac{1}{\sqrt{m}}\E_{C_L}\!\left[|w_L^0(C_L)|^m\right]^{1/m} &\le K.
\end{align*}
\end{assumption}

\begin{assumption}[Diversity]\label{assump:diversity-lowrank}
The support of $\rho^1$ (the measure on $\Omega_1$ for frozen random features $f_1(C_1)$) is $\R^d$ (or dense in $\R^d$). This ensures that the random features $\varphi_1(\ip{f_1(C_1)}{\cdot})$ have dense span in $L^2(\mathcal{P}_X)$.

\textbf{Key difference from full-width case:} Since $f_1(C_1)$ are frozen, this property is maintained trivially throughout training, unlike the full-width case which requires algebraic topology to prove support preservation.
\end{assumption}

\begin{assumption}[Regularity]\label{assump:regularity-lowrank}
There exists $K\ge 1$ such that:
\begin{itemize}
\item For each $\ell\in\{2,\dots,L\}$, $\xi_\ell$ are bounded on compact intervals and piecewise continuous.
\item For each $\ell\in\{2,\dots,L\}$, $\varphi_\ell,\varphi_\ell'$ are $K$-bounded and $K$-Lipschitz; moreover $\varphi_\ell'$ is non-vanishing for $\ell\in\{2,\dots,L-1\}$.
\item For each $\ell\in\{2,\dots,L-1\}$, the mixing kernels $L^{(\ell)}$ satisfy $\sup_{c_\ell,k}|L^{(\ell)}_{c_\ell,k}|\le K$ and $\|L^{(\ell)}\|_{\infty,1} = \sup_{c_\ell}\sum_{k=1}^r |L^{(\ell)}_{c_\ell,k}| \le rK$.
\item $\partial_2\mathcal{L}(\cdot,\cdot)$ is $K$-Lipschitz in the second variable and $K$-bounded.
\item $|X|\le K$ with probability 1.
\item $\norm{f_1(c_1)}\le K$ for all $c_1$.
\end{itemize}
\end{assumption}

\begin{assumption}[Non-degeneracy]\label{assump:nondegeneracy-lowrank}
If the output layer is trained (i.e., $\xi_L(\cdot) \ne 0$), we assume that the initial loss is better than the trivial predictor:
\[
\mathscr{L}(w_1^0, \ldots, w_L^0) < \E_Z[\mathcal{L}(Y, \varphi_L(0))],
\]
where $\varphi_L(0)$ corresponds to the network output when all weights are zero (the trivial predictor $\hat{y} = 0$).

If the output layer is frozen (i.e., $\xi_L(\cdot) = 0$) and the initialization has non-zero mass, then non-degeneracy is automatic and this assumption is not needed.
\end{assumption}

\begin{remark}[On the non-degeneracy assumption]
This assumption ensures that the limit point is non-degenerate (i.e., not all weights are zero). It is used in the global convergence proof to show that $\Pp(\bar{w}_L(C_L) \ne 0) > 0$, which is necessary for the optimality argument.

\textbf{Why this assumption is natural:}
\begin{itemize}
\item For most reasonable initializations (e.g., small random weights, Xavier initialization), the initial network typically performs better than the trivial solution $\hat{y} = 0$, so this condition is satisfied with high probability.
\item The assumption can be easily verified by computing the initial loss value.
\item If this condition fails, one can re-initialize the network or use a different initialization scheme.
\item In the full-width case, the analogous assumption is $\mathscr{L}(w_1^0, w_2^0, w_3^0) < \E_Z[\mathcal{L}(Y, \varphi_3(0))]$ for three-layer networks.
\end{itemize}

\textbf{What happens if the assumption fails:} If $\mathscr{L}(w_1^0, \ldots, w_L^0) \ge \E_Z[\mathcal{L}(Y, \varphi_L(0))]$, then the network starts worse than or equal to the trivial solution. Convergence to a global minimizer is still possible (the global minimizer must be at least as good as the trivial solution), but the non-degeneracy argument needs to be modified. The limit point could potentially be degenerate (all weights zero), which would require a different proof strategy.
\end{remark}

\begin{assumption}[Convergence]\label{assump:convergence-lowrank}
There exist functions $\bar{w}_\ell$ for $\ell\in\{2,\dots,L\}$ such that as $t\to\infty$, there exists a coupling $\pi_t$ such that for each intermediate layer $\ell\in\{2,\dots,L-1\}$:
\[
\int (1+|\bar{w}_L|)\prod_{j=\ell+1}^{L-1}(1+|\bar{w}_j|)\max_{1\le k\le r}|\bar{w}_\ell(c_\ell,k)|\,|w_\ell^*(t,c_\ell',k)-\bar{w}_\ell(c_\ell,k)|\,d\pi_t \to 0,
\]
and for the output layer:
\[
\int (1+|\bar{w}_L|)\,|w_L^*(t,c_L')-\bar{w}_L(c_L)|\,d\pi_t \to 0.
\]
\end{assumption}

\begin{remark}
The convergence assumption is adapted from the full-width case. The key difference is that we have $r$ channels for each intermediate layer, so we take the maximum over $k\in\{1,\dots,r\}$. The product structure $(1+|\bar{w}_L|)\prod_{j=\ell+1}^{L-1}(1+|\bar{w}_j|)$ ensures that the bound remains finite and accounts for the depth $L$.
\end{remark}

\subsection{How We Handle Trainability and Initialization Assumptions}

A critical question is: how does our low-rank framework handle the trainability and initialization assumptions required in the full-width case \cite{nguyen2023rigorousframeworkmeanfield}? We address this systematically:

\paragraph*{The Full-Width Case Requirements (for $L \ge 3$):}

In the full-width case, global convergence for $L \ge 3$ layers requires:

\begin{enumerate}
\item \textbf{Trainability of first layer}: The first layer weights $w_1(t, C_1)$ must be trainable and maintain full support $\text{supp}(w_1(t, C_1)) = \mathbb{R}^d$ throughout training.

\item \textbf{Correlated initialization}: Standard i.i.d. initialization does \textbf{not} work for $L \ge 3$. Instead, one needs a \textbf{correlated initialization} scheme where the joint distribution $(w_1^0(C_1), w_2^0(C_1, \cdot))$ has full support in $\mathbb{R}^d \times L^2(P_2)$.

\item \textbf{Bidirectional diversity}: For $L \ge 3$, the initialization must satisfy:
\[
\text{supp}(w_1^0(C_1), w_2^0(C_1, \cdot)) = \mathbb{R}^d \times L^2(P_2),
\]
and for intermediate layers $i = 2, \ldots, L-1$:
\[
\text{supp}(w_i^0(\cdot, C_i), w_{i+1}^0(C_i, \cdot)) = L^2(P_{i-1}) \times L^2(P_{i+1}).
\]

\item \textbf{Function-dependent design}: The correlated initialization scheme is \textbf{function-dependent}---the correlation structure must be chosen based on the target function to approximate, making it non-agnostic.
\end{enumerate}

\paragraph*{Why These Requirements Are Needed in Full-Width Case:}

The fundamental issue is that with i.i.d. initialization for $L \ge 3$, neurons at intermediate layers \textbf{collapse to the same function} (Corollary 29 of \cite{nguyen2023rigorousframeworkmeanfield}). This means:
\[
w_i(t, c_{i-1}, c_i) \approx w_i(t, c_{i-1}', c_i') \quad \text{(all neurons learn the same function)}
\]
for intermediate layers, breaking universal approximation. The correlated initialization with bidirectional diversity is designed to prevent this collapse.

\paragraph*{How Our Low-Rank Framework Eliminates These Requirements:}

\begin{enumerate}
\item \textbf{No trainability assumption for first layer}: We use \textbf{frozen random features} $f_1(C_1)$ in the first layer. Since $f_1(C_1)$ is fixed and not trained, we don't need to prove that it maintains full support---it trivially does because it never changes.

\item \textbf{Standard i.i.d. initialization suffices}: We can initialize $w_1^0(C_1, k)$ and $w_2^0(C_2)$ independently from standard distributions (e.g., Gaussian), without any correlation structure. The frozen random features provide the necessary diversity.

\item \textbf{No bidirectional diversity needed}: The frozen random features $f_1(C_1)$ automatically provide full support $\text{supp}(f_1(C_1)) = \mathbb{R}^d$ throughout training, which is sufficient for universal approximation (Theorem~\ref{thm:relu-dense-span}). We don't need the complex bidirectional diversity structure.

\item \textbf{Function-agnostic}: The initialization is completely independent of the target function. We don't need to design a correlation structure based on what we're trying to approximate.
\end{enumerate}

\paragraph*{The Key Mechanism:}

The frozen first layer solves the collapse problem at its source:
\begin{itemize}
\item In full-width networks, if the first layer is trainable with i.i.d. initialization, neurons can collapse during training, breaking universal approximation.
\item In our framework, the first layer is frozen, so it cannot collapse. The frozen random features maintain full support $\text{supp}(f_1(C_1)) = \mathbb{R}^d$ trivially.
\item This provides a stable, diverse basis that enables universal approximation without needing correlated initialization or bidirectional diversity.
\end{itemize}

\paragraph*{Comparison Summary:}

\begin{center}
\begin{tabular}{|l|l|l|}
\hline
\textbf{Requirement} & \textbf{Full-Width ($L \ge 3$)} & \textbf{Low-Rank (any $L$)} \\
\hline
First layer trainability & Required (with support preservation) & Not needed (frozen) \\
Initialization & Correlated (function-dependent) & Standard i.i.d. \\
Bidirectional diversity & Required & Not needed \\
Function-agnostic & No & Yes \\
Depth limit & $L = 3$ with i.i.d. & Arbitrary $L$ \\
\hline
\end{tabular}
\end{center}

\textbf{Conclusion:} Our low-rank framework eliminates the need for trainability assumptions on the first layer and correlated initialization schemes by using frozen random features. This makes the framework simpler, more practical, and applicable to arbitrary depth $L$ with standard initialization.

\section{Existence and Uniqueness of the Solution of the MF ODEs}
\label{sec:Existence-MF}

We establish well-posedness of the mean-field ODE system \eqref{eq:ode-w1}--\eqref{eq:ode-wL} by adapting the proof from \cite{nguyen2023rigorousframeworkmeanfield} to account for the low-rank structure.

\begin{theorem}[Well-posedness of mean-field ODEs]
\label{thm:wellposed}
Under Assumptions~\ref{assump:init-lowrank}, \ref{assump:regularity-lowrank}, and \ref{assump:diversity-lowrank}, there exists a unique solution to the mean-field ODE system \eqref{eq:ode-w1}--\eqref{eq:ode-wL} on $t\in[0,\infty)$.
\end{theorem}

\begin{proof}[Proof sketch]
The proof reveals how the low-rank structure creates norms that account for the $r$ independent channels. We use Picard iteration (following \cite{nguyen2023rigorousframeworkmeanfield}), but the key difference is that the low-rank mixing matrix $L$ requires function spaces equipped with norms that account for the $r$ independent channels. We establish a priori bounds showing that solutions remain bounded in these multi-channel spaces. The solution operator $F$ maps these spaces into themselves and is contractive. Applying the Banach fixed point theorem yields existence and uniqueness. The crucial insight is that the multiplicative factors of $r$ in bounds reflect the channel specialization structure: each of the $r$ channels contributes independently. The detailed proof, including all technical lemmas and norm definitions, is deferred to Appendix~\ref{app:wellposed-proof}.
\end{proof}

\section{Main Result: Connection between Neural Network and MF Limit}
\label{sec:Main-result}

This section establishes the quantitative approximation result, showing that finite-width networks converge to the mean-field limit as width increases.

\begin{theorem}[Quantitative approximation for low-rank networks]
\label{thm:quantitative-lowrank}
Under the regularity assumptions, for any $T>0$ and $\delta>0$, with probability at least $1-\delta$, the distance between the finite-width network and the mean-field limit satisfies:
\[
\mathscr{D}_T(W, \mathbf{W}) \le e^{K_T(1+rK)^{L-2}}\left(\frac{1}{\sqrt{n_{\min}}} + \sqrt{\epsilon}\right)\log^{1/2}\left(\frac{3(T+1)n_{\max}^2r^{L-2}}{\delta} + e\right),
\]
where $n_{\min} = \min\{n_1, \ldots, n_{L-1}\}$ and $n_{\max} = \max\{n_1, \ldots, n_{L-1}\}$.
\end{theorem}

\begin{remark}[Generalization to $L$ layers]
\label{rem:generalization-L}
The proof strategy generalizes naturally to arbitrary depth $L$ because \textbf{the same three-step pattern repeats at each intermediate layer}:

\begin{enumerate}
\item \textbf{Forward decomposition}: For each $H_i$ with $i = 2, \ldots, L$, we decompose the difference using the same structure with factors of $rK$ from $\|L^{(i)}\|_{\infty,1} \le rK$.

\item \textbf{Backward decomposition}: For each backpropagated signal $B_k^{(i)}$ with $i = 2, \ldots, L-1$, we decompose similarly.

\item \textbf{Concentration and Gronwall}: The concentration bounds and Gronwall argument proceed identically at each layer, with union bounds over $r$ channels.
\end{enumerate}

\textbf{Key insight:} With bounded mixing matrices $\|L^{(i)}\|_{\infty,1} \le rK$, bounded activations, and bounded data $|X| \le K$, the exponential constant accumulates multiplicatively as $(1+rK)^{L-2}$ rather than exponentially in $L$. The frozen first layer eliminates recursion at the base case, ensuring the pattern starts cleanly.
\end{remark}

\section{Convergence to Global Optimum: Low-Rank Networks with I.i.d. Initialization}
\label{sec:global_convergence_iid}

In this section, we prove global convergence guarantees for low-rank random feature networks of arbitrary depth $L \ge 2$ with standard i.i.d. initialization. A key element here is a certain universal approximation property that holds at \textit{any} finite training time, which is automatically maintained by frozen random features.

\subsection{Main Result}

\begin{theorem}[Global convergence for low-rank networks of arbitrary depth]
\label{thm:global-convergence-lowrank}
Under Assumptions~\ref{assump:init-lowrank}--\ref{assump:convergence-lowrank}, if the loss $\mathcal{L}$ is convex in its second variable, then the limit point $(\bar{w}_1,\ldots,\bar{w}_L)$ is a global minimizer of the population loss:
\[
\mathscr{L}(\bar{w}_1,\ldots,\bar{w}_L) = \E_Z\big[\mathcal{L}\big(Y,\hat y(X;\bar{w}_1,\ldots,\bar{w}_L)\big)\big].
\]
For any depth $L \ge 2$, this holds with standard i.i.d. initialization, unlike full-width networks which require correlated initialization for $L \ge 3$.
\end{theorem}

\subsection{High-Level Proof Strategy}

The proof follows the same structure as the full-width case but with key simplifications:

\begin{enumerate}
\item \textbf{Convergence of backpropagated signal}: Show that the backpropagated signals $B_k^{(i)}$ converge to zero at the limit point.

\item \textbf{Universal approximation property}: Use the fact that frozen random features $f_1(C_1)$ maintain full support $\text{supp}(f_1(C_1)) = \R^d$ throughout training, which implies dense span in $L^2(\mathcal{P}_X)$ (see Theorem~\ref{thm:relu-dense-span}).

\item \textbf{Orthogonality and optimality}: Show that the gradient is orthogonal to all functions in the dense span, implying global optimality.

\item \textbf{Non-degeneracy}: Ensure the limit point is non-degenerate.

\item \textbf{Global optimality}: Use convexity of the loss to conclude global optimality.
\end{enumerate}

\subsection{Detailed Proof}

\begin{proof}[Proof of Theorem~\ref{thm:global-convergence-lowrank}]

\textbf{Step 1: Convergence of the backpropagated signal}

By the convergence assumption (Assumption~\ref{assump:convergence-lowrank}), we have that for any $\epsilon>0$, there exists $T(\epsilon)$ such that for all $t\ge T(\epsilon)$ and almost surely:
\[
\epsilon \ge \left|\E_{Z}\!\left[d_L\big(Z;W^*(t)\big)\,\varphi_1\!\left(\ip{f_1(C_1)}{X}\right)\,B_k^{(2)*}(t;X)\right]\right|,
\]
where $B_k^{(2)*}(t;X)$ is the backpropagated signal at the solution $W^*(t)$.

\textbf{Step 2: Universal approximation property}

Since $f_1(C_1)$ are frozen random features with full support in $\R^d$ (by Assumption~\ref{assump:diversity-lowrank}), and we have $r$ channels, the set $\{\varphi_1(\ip{f_1(c_1)}{\cdot}): c_1\in\Omega_1, k\in\{1,\dots,r\}\}$ has dense span in $L^2(\mathcal{P}_X)$ by Theorem~\ref{thm:relu-dense-span}.

\textbf{Step 3: Orthogonality and optimality}

Using the dense span property and the convergence of the backpropagated signal, we obtain that for $\mathcal{P}_X$-almost every $x$:
\[
\E_{Z}\!\left[\partial_2\mathcal{L}\big(Y,\hat y(X;\bar{W})\big)\middle|X=x\right] = 0.
\]

\textbf{Step 4: Non-degeneracy}

We ensure that the limit point is non-degenerate. This follows from the assumption that the initial loss is better than the trivial predictor, combined with the gradient flow property (loss decreases along trajectories).

\textbf{Step 5: Global optimality}

Since $\mathcal{L}$ is convex in its second variable, for any measurable function $\tilde{y}(x)$:
\[
\E_Z\big[\mathcal{L}\big(Y,\tilde{y}(X)\big)\big] \ge \mathscr{L}(\bar{w}_1,\ldots,\bar{w}_L),
\]
i.e., $(\bar{w}_1,\ldots,\bar{w}_L)$ is a global minimizer.

\end{proof}

\subsection{Key Differences from Full-Width Case}

\begin{enumerate}
\item \textbf{Simplification: No algebraic topology needed}: Since we do not train $\rho^1$ (the first layer uses frozen random features), we avoid the complex arguments needed to show that $\text{Law}(w_1^*(t,C_1))$ has full support. Instead, we rely on the fact that $f_1(C_1)$ are frozen random features with full support, which directly gives the dense span property.

\item \textbf{Low-rank structure}: The pre-activations $H_i$ are sums over $r$ channels, introducing factors of $\|L^{(i)}\|_{\infty,1}\le rK$ in bounds.

\item \textbf{Convergence assumption}: Instead of bounding differences for a single $w_1$, we take the maximum over $r$ channels: $\max_{1\le k\le r}|w_1^*(t,c_1',k)-\bar{w}_1(c_1,k)|$.

\item \textbf{Arbitrary depth}: The proof works for any $L \ge 2$, not just $L=3$, because the frozen first layer provides a stable base and the low-rank structure ensures polynomial (not exponential) growth of constants.
\end{enumerate}

\section{Universal Approximation Property: ReLU Dense Span}
\label{sec:relu-dense-span}

This section provides a rigorous proof that frozen random features with full support give dense span in $L^2(\mathcal{P}_X)$, adapting the approach from the full-width case.

\begin{theorem}[Dense span for frozen ReLU random features]
\label{thm:relu-dense-span}
Let $\varphi_1$ be the ReLU activation function, and let $f_1(C_1)$ be frozen random feature vectors with $\text{supp}(f_1(C_1)) = \mathbb{R}^d$. Then the function class:
\[
\mathcal{F} = \{\varphi_1(\langle f_1(c_1), \cdot\rangle) : c_1 \in \Omega_1\}
\]
has dense span in $L^2(\mathcal{P}_X)$ (with respect to the data distribution $\mathcal{P}_X$).
\end{theorem}

\begin{proof}[Proof sketch]
The proof proceeds in three steps:

\begin{enumerate}
\item \textbf{Standard universal approximation result}: Use Leshno et al. (1993) to show that $\{\varphi_1(\langle u, \cdot\rangle) : u \in \mathbb{R}^d\}$ has dense span in $L^2(\mathcal{P}_X)$.

\item \textbf{Full support gives dense parameterization}: Since $\text{supp}(f_1(C_1)) = \mathbb{R}^d$, the set $\{f_1(c_1) : c_1 \in \Omega_1\}$ is dense in $\mathbb{R}^d$.

\item \textbf{Continuity argument}: Use continuity of $u \mapsto \varphi_1(\langle u, \cdot\rangle)$ in $L^2(\mathcal{P}_X)$ to show that dense parameterization implies dense span.
\end{enumerate}

The detailed proof is provided in Appendix~\ref{app:relu-dense-span-proof}.
\end{proof}

\section{Channel Specialization Mechanism}
\label{sec:channel-specialization}

The low-rank structure enables channel specialization where different channels capture different frequency regimes. The following theorem provides a rigorous characterization of how channels establish and maintain dominance.

\subsection{Why Channels Cannot Learn the Same Function: The Role of $L_{c_2,k}$}

A natural question is: why don't all $r$ channels simply learn the same function $f_1 = f_2 = \cdots = f_r$? This would make the low-rank structure redundant. The answer lies in how the mixing matrix $L_{c_2,k}$ creates \textbf{channel-specific backpropagated signals} that force different gradients for different channels.

\begin{proposition}[Mixing matrix prevents channel collapse]
\label{prop:no-collapse}
Suppose that at some time $t$, all channels have learned the same function: $f_1(t,\cdot) = f_2(t,\cdot) = \cdots = f_r(t,\cdot) = f(t,\cdot)$. Then the backpropagated signals are:
\[
B_k(t; X) = \E_{C_2}\!\left[L_{C_2,k}\,\varphi_2'\!\Big(H_2(t, C_2; X)\Big)\,w_2(t, C_2)\right],
\]
where $H_2(t, c_2; X) = \sum_{j=1}^r L_{c_2,j} f(t, X) = f(t, X) \sum_{j=1}^r L_{c_2,j}$.

If the mixing matrix $L$ is \textbf{non-isotropic} (i.e., the columns $L_{\cdot,k}$ are not all proportional to each other), then the backpropagated signals $B_k$ are \textbf{not all equal}, even when $f_k = f$ for all $k$. This creates different gradients $\partial_t w_1(t, c_1, k) \propto \E_Z[d_L \varphi_1 B_k]$, which forces the channels to diverge and learn different functions.
\end{proposition}

\begin{proof}[Proof sketch]
The key observation is that even if $f_1 = f_2 = \cdots = f_r = f$, we have:
\[
B_k(t; X) = \E_{C_2}\!\left[L_{C_2,k}\,\varphi_2'\!\Big(f(t, X) \sum_{j=1}^r L_{c_2,j}\Big)\,w_2(t, C_2)\right].
\]

Since $\varphi_2'$ is non-constant (e.g., ReLU gives $\varphi_2'(H_2) = \mathbf{1}\{H_2 > 0\}$), the gate $\varphi_2'(f(t, X) \sum_{j=1}^r L_{c_2,j})$ depends on the \textbf{sum} $\sum_{j=1}^r L_{c_2,j}$, which is the same for all $k$. However, the \textbf{coefficient} $L_{C_2,k}$ in front of the gate is different for each channel $k$.

If the columns $L_{\cdot,k}$ are not all proportional (non-isotropic), then for some neurons $c_2$, we have $L_{c_2,1} \ne L_{c_2,2}$ (up to a constant factor). This means:
\[
B_1(t; X) = \E_{C_2}\!\left[L_{C_2,1}\,\varphi_2'(\cdots)\,w_2\right] \ne \E_{C_2}\!\left[L_{C_2,2}\,\varphi_2'(\cdots)\,w_2\right] = B_2(t; X),
\]
even when $f_1 = f_2 = f$.

Since $\partial_t w_1(t, c_1, k) = -\xi_1(t) \E_Z[d_L \varphi_1 B_k]$, different $B_k$ create different gradients, forcing the channels to learn different functions. This prevents collapse.
\end{proof}

\textbf{The mechanism in detail:}

\begin{enumerate}
\item \textbf{Channel-specific mixing coefficients}: Each channel $k$ has its own column $L_{\cdot,k}$ of the mixing matrix. When computing $B_k$, we weight by $L_{C_2,k}$, which is different for each $k$.

\item \textbf{ReLU gating creates competition}: For ReLU, $\varphi_2'(H_2) = \mathbf{1}\{H_2 > 0\}$. The gate depends on $H_2 = \sum_{j=1}^r L_{c_2,j} f_j$. Even if all $f_j = f$, the gate $\mathbf{1}\{f \sum_j L_{c_2,j} > 0\}$ is the same for all channels, but the coefficient $L_{c_2,k}$ in $B_k$ is different.

\item \textbf{Non-isotropic structure breaks symmetry}: If $L$ is isotropic (all columns proportional), then $B_k = c(t) f_k$ for all $k$, and channels can collapse. But if $L$ is non-isotropic (e.g., deterministic grid, structured initialization), then $B_k$ are not proportional to $f_k$, creating different gradients.

\item \textbf{Self-reinforcing divergence}: Once channels start to diverge (due to different $B_k$), the log-ratio growth mechanism (Theorem~\ref{thm:r2-logratio-criterion-main}) amplifies the differences, preventing collapse.
\end{enumerate}

\textbf{Counterexample: Isotropic mixing allows collapse}

If $L_{C_2} \sim \mathcal{N}(0, I_r)$ (isotropic Gaussian), then in expectation:
\[
\E[L_{C_2,k} | f_1, \ldots, f_r] \propto f_k,
\]
which makes $B_k \propto f_k$. If all channels start with $f_1 = f_2 = \cdots = f_r$, then all $B_k$ are equal, and channels can remain identical. This is why \textbf{non-isotropic mixing is crucial} for preventing collapse.

\textbf{Key insight:} The mixing matrix $L_{c_2,k}$ acts as a \textbf{differentiation mechanism}. By giving each channel $k$ a different column $L_{\cdot,k}$, we ensure that even if channels start identically, they receive different backpropagated signals $B_k$, which forces them to learn different functions. The non-isotropic structure of $L$ is what breaks the symmetry and enables specialization.

\begin{theorem}[Channels learn independent spikes]
\label{thm:r2-logratio-criterion-main}
Fix the spike location $x_0$ and assume the local reduction holds there:
\[
\partial_t f_k(t,x_0)=-\xi_1(t)\,K_{\mu_0}(x_0,x_0)\,d_0(t)\,B_{k,1}(t),\qquad k=1,2,
\]
where $d_0(t)$ is the residual signal at $x_0$ and $B_{k,1}(t)=B_k(t;x_0)$. Assume there exists a constant $\rho\in[0,1)$ such that on a time interval $I$,
\begin{equation}
\label{eq:r2-stability-assumption-main}
\big|B_{2,1}(t)\big|\le \rho\,\frac{|f_2(t,x_0)|+\varepsilon}{|f_1(t,x_0)|+\varepsilon}\,\big|B_{1,1}(t)\big|
\qquad\text{for all }t\in I.
\end{equation}
If moreover $-d_0(t)\ge 0$ and $B_{1,1}(t)$ has the same sign as $f_1(t,x_0)$ on $I$, then the log-ratio
\[
R_{12}(t,x_0)=\log\frac{|f_1(t,x_0)|+\varepsilon}{|f_2(t,x_0)|+\varepsilon}
\]
satisfies $\partial_t R_{12}(t,x_0)\ge (1-\rho)\,\xi_1(t)\,K_{\mu_0}(x_0,x_0)\,(-d_0(t))\,\frac{|B_{1,1}(t)|}{|f_1(t,x_0)|+\varepsilon}\ge 0$ for a.e.\ $t\in I$. In particular, any strict dominance at $t=t_0\in I$ cannot be lost on $I$ and is amplified whenever $|B_{1,1}|$ is not too small.
\end{theorem}

\begin{proof}[Proof sketch]
The proof is deferred to Appendix~\ref{app:r2-logratio-proof}. The key idea is that the stability assumption \eqref{eq:r2-stability-assumption-main} ensures that the backpropagated signal for the second channel is bounded relative to the first channel, which preserves and amplifies channel dominance during training.

\textbf{Dominance in the integral $\partial_t f_k$:} The evolution equation $\partial_t f_k(t,x)=-\xi_1(t)\,\E_{Z=(X,Y)}[d_L(Z;W(t))\,B_k(t;X)\,K_{\mu_0}(x,X)]$ is an integral over the data distribution. When channel 1 dominates at $x_0$ (i.e., $|f_1(t,x_0)|>|f_2(t,x_0)|$), the sign-coherence condition ($B_{1,1}$ has the same sign as $f_1$) ensures that the integral contribution from channel 1 is larger in magnitude than that from channel 2. This creates a self-reinforcing effect: dominance leads to larger backpropagated signals $B_k$, which make the integral $\partial_t f_k$ larger, further amplifying the dominance.

\textbf{$L_{c_2}$ sign correlation intuition:} The backpropagated signal $B_k(t;X)=\E_{C_2}[L_{C_2,k}\,w_2(t,C_2)\,\1\{H_2(t,C_2;X)>0\}]$ is a gated mixture of the $k$th column of $L$. When $L_{c_2,k}$ has consistent sign (or controlled sign changes) and $w_2(t,c_2)$ is aligned, increasing $f_1(t,\cdot)$ increases $H_2=\sum_{j=1}^r L_{c_2,j}f_j$ on the subset where $L_{c_2,1}>0$, which increases the gate $\1\{H_2>0\}$ and hence increases $B_1(t;X)$. Since $\partial_t f_1$ is proportional to $B_1$ through the integral, this creates a positive feedback loop: larger $f_1$ $\to$ larger $H_2$ $\to$ larger gate $\to$ larger $B_1$ $\to$ larger $\partial_t f_1$ $\to$ even larger $f_1$.
\end{proof}

\section{Appendices}

\subsection{Proof of Well-Posedness}
\label{app:wellposed-proof}

We establish well-posedness of the mean-field ODE system \eqref{eq:ode-w1}--\eqref{eq:ode-wL} by adapting the proof from \cite{nguyen2023rigorousframeworkmeanfield} to account for the low-rank structure.

\subsubsection{Norms and Spaces}

We equip the mean-field parameters with several norms. For the low-rank case with $L$ layers and $r$ channels, we define:

\begin{align*}
\interleave w_1\interleave_t &= \max_{1\le k\le r}\E_{C_1}\!\left[\sup_{s\le t}|w_1(s,C_1,k)|^{50}\right]^{1/50},\\
\interleave w_i\interleave_t &= \max_{1\le k\le r}\E_{C_i}\!\left[\sup_{s\le t}|w_i(s,C_i,k)|^{50}\right]^{1/50}, \quad i=2,\ldots,L-1,\\
\interleave w_L\interleave_t &= \E_{C_L}\!\left[\sup_{s\le t}|w_L(s,C_L)|^{50}\right]^{1/50},\\
\interleave W\interleave_t &= \max\big(\interleave w_1\interleave_t, \max_{2\le i\le L-1}\interleave w_i\interleave_t, \interleave w_L\interleave_t\big).
\end{align*}

We also define the $L^2$-type norms:
\begin{align*}
\|w_1\|_t &= \max_{1\le k\le r}\E_{C_1}\!\left[\sup_{s\le t}|w_1(s,C_1,k)|^2\right]^{1/2},\\
\|w_i\|_t &= \max_{1\le k\le r}\E_{C_i}\!\left[\sup_{s\le t}|w_i(s,C_i,k)|^2\right]^{1/2}, \quad i=2,\ldots,L-1,\\
\|w_L\|_t &= \E_{C_L}\!\left[\sup_{s\le t}|w_L(s,C_L)|^2\right]^{1/2},\\
\|W\|_t &= \max\big(\|w_1\|_t, \max_{2\le i\le L-1}\|w_i\|_t, \|w_L\|_t\big).
\end{align*}

For the $\psi_2$-type sub-Gaussian norms:
\begin{align*}
\llbracket w_1\rrbracket_{\psi,t} &= \sqrt{50}\,\sup_{m\ge 1}\frac{1}{\sqrt{m}}\max_{1\le k\le r}\E_{C_1}\!\left[\sup_{s\le t}|w_1(s,C_1,k)|^m\right]^{1/m},\\
\llbracket w_i\rrbracket_{\psi,t} &= \sqrt{50}\,\sup_{m\ge 1}\frac{1}{\sqrt{m}}\max_{1\le k\le r}\E_{C_i}\!\left[\sup_{s\le t}|w_i(s,C_i,k)|^m\right]^{1/m}, \quad i=2,\ldots,L-1,\\
\llbracket w_L\rrbracket_{\psi,t} &= \sqrt{50}\,\sup_{m\ge 1}\frac{1}{\sqrt{m}}\E_{C_L}\!\left[\sup_{s\le t}|w_L(s,C_L)|^m\right]^{1/m},\\
\llbracket W\rrbracket_{\psi,t} &= \max\big(\llbracket w_1\rrbracket_{\psi,t}, \max_{2\le i\le L-1}\llbracket w_i\rrbracket_{\psi,t}, \llbracket w_L\rrbracket_{\psi,t}\big).
\end{align*}

The factor $\sqrt{50}$ ensures that $\llbracket W\rrbracket_{\psi,t}\ge \interleave W\interleave_t$ and hence $\llbracket W\rrbracket_{\psi,t}\ge \|W\|_t$.

For convenience, we define the random variables:
\begin{align*}
\mathsf{max}_t^w(W) &= \max_{1\le k\le r}\sup_{s\le t}|w_1(s,C_1,k)|,\\
\mathsf{max}_t^{w_i}(W) &= \max_{1\le k\le r}\sup_{s\le t}|w_i(s,C_i,k)|, \quad i=2,\ldots,L-1,\\
\mathsf{max}_t^{w_L}(W) &= \sup_{s\le t}|w_L(s,C_L)|.
\end{align*}

For two sets of mean-field parameters $W$ and $W'$, we define the distance:
\begin{align}
\|W-W'\|_t &= \max\big(\|w_1-w_1'\|_t, \max_{2\le i\le L-1}\|w_i-w_i'\|_t, \|w_L-w_L'\|_t\big),\label{eq:def-MF-distance-app}\\
\|w_1-w_1'\|_t &= \max_{1\le k\le r}\E_{C_1}\!\left[\sup_{s\le t}|w_1(s,C_1,k)-w_1'(s,C_1,k)|^2\right]^{1/2},\nonumber\\
\|w_i-w_i'\|_t &= \max_{1\le k\le r}\E_{C_i}\!\left[\sup_{s\le t}|w_i(s,C_i,k)-w_i'(s,C_i,k)|^2\right]^{1/2}, \quad i=2,\ldots,L-1,\nonumber\\
\|w_L-w_L'\|_t &= \E_{C_L}\!\left[\sup_{s\le t}|w_L(s,C_L)-w_L'(s,C_L)|^2\right]^{1/2}.\nonumber
\end{align}

\subsubsection{A Priori Bounds}

\begin{lemma}[Weight bounds]
\label{lem:bounds-a-priori}
Under Assumptions~\ref{assump:regularity-lowrank} and \ref{assump:init-lowrank}, given an initialization $W(0)$, a solution $W$ to the mean-field ODEs, if it exists, must satisfy that for any $t\in[0,\infty)$:
\[
\interleave W\interleave_t \le K_0(t),
\]
where $K_0(t)$ is a non-decreasing function of the form
\[
K_0(t) = (1+rK)^{1/2}\,K^{\kappa}(1+t^{\kappa})(1+\interleave W\interleave_0^{\kappa}),
\]
for some constant $\kappa>0$ depending on $K$ and $r$.

A similar result holds for the $\psi_2$ norm. Under the same assumptions, for any $t\in[0,\infty)$, there exists $K_0(t)\ge 1$ of the form
\[
K_0(t) = (1+rK)^{1/2}\,K^{\kappa}(1+t^{\kappa})(1+\llbracket W\rrbracket_{\psi,0}^{\kappa}),
\]
such that a solution $W$, if it exists, must satisfy $\interleave W\interleave_t \le \llbracket W\rrbracket_{\psi,t} \le K_0(t)$ for any $t\in[0,\infty)$. Furthermore, by assuming $\llbracket W\rrbracket_{\psi,0}<\infty$, for any $B\ge 0$:
\[
\Pp\!\left(\max\big(\mathsf{max}_t^w(W), \mathsf{max}_t^{w_2}(W)\big)\ge K_0(t)B\right) \le C e^{-K_1 B^2},
\]
for some universal constants $C, K_1>0$, where $\mathsf{max}_t^w(W) = \max_{1\le k\le r}\sup_{s\le t}|w_1(s,C_1,k)|$ and $\mathsf{max}_t^{w_2}(W) = \sup_{s\le t}|w_2(s,C_2)|$.
\end{lemma}

\begin{proof}
The proof follows from Grönwall-type arguments applied to the ODEs \eqref{eq:ode-w1}--\eqref{eq:ode-wL}. The key steps are:
\begin{enumerate}
\item Use the boundedness and Lipschitz properties of $\varphi_i$, $\varphi_i'$ for $i=1,\ldots,L$, and $\partial_2\mathcal{L}$.

\item Control $H_i$ for $i=2,\ldots,L$ using the low-rank structure:
\begin{align*}
|H_2| &\le \|L^{(2)}\|_{\infty,1}\max_{1\le k\le r}|m_k^{(2)}|\le rK\max_{1\le k\le r}\E_{C_1}[|w_1(C_1,k)|],\\
|H_i| &\le \|L^{(i)}\|_{\infty,1}\max_{1\le k\le r}|m_k^{(i)}|\le rK\max_{1\le k\le r}\E_{C_{i-1}}[|w_{i-1}(C_{i-1},k)|], \quad i=3,\ldots,L.
\end{align*}

\item Control $B_k^{(i)}$ for $i=2,\ldots,L-1$ using the structure:
\begin{align*}
|B_k^{(2)}| &\le K\E_{C_2}[|L^{(2)}_{C_2,k}|\,|w_2|]\le rK^2\E_{C_2}[|w_2|],\\
|B_k^{(i)}| &\le K\E_{C_i}[|L^{(i)}_{C_i,k}|\,|w_i|]\le rK^2\E_{C_i}[|w_i|], \quad i=3,\ldots,L-1.
\end{align*}

\item Apply Minkowski's inequality and Grönwall's lemma to obtain polynomial growth in $t$. The key is that each layer contributes a factor $(1+rK)$ to the Gronwall constant, giving $(1+rK)^{L-2}$ total, which is polynomial in $L$ (not exponential).

\item The sub-Gaussian tail bound follows from the $\psi_2$ control and a union bound over $r$ channels and $L$ layers.
\end{enumerate}
The constant $\kappa$ depends on $K$, $r$, and $L$ through the factors $\|L^{(i)}\|_{\infty,1}\le rK$ for $i=2,\ldots,L-1$, but remains polynomial rather than exponential in $r$ and $L$ because the factors multiply (not compound exponentially).
\end{proof}

\subsubsection{Solution Operator and Fixed Point Formulation}

Denote by $\mathfrak{W}_T$ the space of mean-field parameters $W$ such that $\|W\|_T<\infty$. Given a terminal time $T\ge 0$ and an initialization $W(0)$, we define the mapping $F$ that associates $W'\in\mathfrak{W}_T$ with:
\begin{align*}
F(W')(t) &= \big\{F_1(W')(t,\cdot,\cdot), F_2(W')(t,\cdot), \ldots, F_L(W')(t,\cdot)\big\},
\end{align*}
where for $i=1,\ldots,L-1$ and $k=1,\ldots,r$:
\begin{align*}
F_1(W')(t,c_1,k) &= w_1(0,c_1,k) - \int_0^t \xi_1(s)\,\E_{Z}\!\left[d_L\big(Z;W'(s)\big)\,\varphi_1\!\left(\ip{f_1(c_1)}{X}\right)\,B_k^{(2)}\big(s;X,W'(s)\big)\right]ds,\\
F_i(W')(t,c_i,k) &= w_i(0,c_i,k) - \int_0^t \xi_i(s)\,\E_{Z}\!\left[d_L\big(Z;W'(s)\big)\,\varphi_i\!\Big(H_i\big(s,c_i;X,W'(s)\big)\Big)\,B_k^{(i+1)}\big(s;X,W'(s)\big)\right]ds,
\end{align*}
and for the output layer:
\begin{align*}
F_L(W')(t,c_L) &= w_L(0,c_L) - \int_0^t \xi_L(s)\,\E_{Z}\!\left[d_L\big(Z;W'(s)\big)\,\varphi_L\!\Big(H_L\big(s,c_L;X,W'(s)\big)\Big)\right]ds.
\end{align*}

Observe that at initialization $F(W')(0,\cdot,\cdot)=W(0)$, whereas the quantities in the time integrals are computed with respect to $W'$. A solution to the mean-field ODEs on $[0,T]$ is an element of $\mathfrak{W}_T$ satisfying $F(W)=W$. We say that $W$ is a solution on $[0,\infty)$ if its restriction to $[0,T]$ is a solution on $[0,T]$ for all $T>0$.

These a priori bounds lead us to consider the following spaces, given an initialization $W(0)$ and an arbitrary terminal time $T>0$:

\begin{itemize}
\item The space $\mathcal{W}_T$ of mean-field parameters $W'=\{W'(t)\}_{t\le T}$ such that $\interleave W'\interleave_T \le K_0(T)$.

\item The space $\mathcal{W}_T^0\subset\mathcal{W}_T$ of mean-field parameters $W'\in\mathcal{W}_T$ such that:
\begin{align*}
\llbracket W'\rrbracket_{\psi,T} &\le K_0(T),\\
\Pp\!\left(\max\big(\mathsf{max}_T^w(W'), \max_{2\le i\le L-1}\mathsf{max}_T^{w_i}(W'), \mathsf{max}_T^{w_L}(W')\big)\ge K_0(T)B\right) &\le C e^{-K_1 B^2} \quad \forall B\ge 0,
\end{align*}
and $W'(0)=W(0)$ (so all elements in $\mathcal{W}_T^0$ share the same initialization).
\end{itemize}

We equip these spaces with the metric $(W',W'')\mapsto \|W'-W''\|_T$. By Lemma~\ref{lem:bounds-a-priori}, any solution $W$ to the mean-field ODEs, if it exists, must belong to $\mathcal{W}_T^0$.

\begin{lemma}[Solution operator maps bounded sets to bounded sets]
\label{lem:invariance-F}
Under Assumptions~\ref{assump:regularity-lowrank} and \ref{assump:init-lowrank}, for any $W'\in\mathcal{W}_T^0$, we have $F(W')\in\mathcal{W}_T^0$.
\end{lemma}

\begin{proof}
The proof follows the same argument as Lemma~\ref{lem:bounds-a-priori}, using the integral form defining $F$ and the bounded/Lipschitz properties of the drifts. The key is that applying bounded/Lipschitz drifts to $W'$ preserves the moment and tail bounds with constants controlled by $K_0(T)$.
\end{proof}

\begin{lemma}[Solution operator is contractive]
\label{lem:difference-F}
For a given $B\ge 0$, consider two collections of mean-field parameters $W',W''\in\mathcal{W}_T$ such that:
\begin{align*}
\Pp\!\left(\max\big(\mathsf{max}_T^w(W'), \max_{2\le i\le L-1}\mathsf{max}_T^{w_i}(W'), \mathsf{max}_T^{w_L}(W')\big)\ge K_0(T)B\right) &\le C e^{-K_1 B^2},\\
\Pp\!\left(\max\big(\mathsf{max}_T^w(W''), \max_{2\le i\le L-1}\mathsf{max}_T^{w_i}(W''), \mathsf{max}_T^{w_L}(W'')\big)\ge K_0(T)B\right) &\le C e^{-K_1 B^2}.
\end{align*}
Under Assumptions~\ref{assump:regularity-lowrank}--\ref{assump:init-lowrank}, for any $t\le T$:
\[
\|F(W')-F(W'')\|_t \le (KK_0(T))^4(1+rK)^{L-2} \int_0^t \Big((1+B)\|W'-W''\|_s + \sqrt{2}e^{-K_1 B^2/2}\Big)ds,
\]
where the constant depends on $K$, $r$, and $L$ through $\|L^{(i)}\|_{\infty,1}\le rK$ for $i=2,\ldots,L-1$.
\end{lemma}

\begin{proof}
The proof uses a good/bad event decomposition:
\begin{enumerate}
\item On the good event $\{\max(\mathsf{max}_T^w, \max_{2\le i\le L-1}\mathsf{max}_T^{w_i}, \mathsf{max}_T^{w_L})\le K_0(T)B\}$, the drifts are Lipschitz in $W$ with constant $(1+B)K_0(T)$.

\item For $H_2$, the difference is controlled using the low-rank structure:
\[
|H_2(W')-H_2(W'')| \le \|L^{(2)}\|_{\infty,1}\max_{1\le k\le r}\E_{C_1}[|w_1'(C_1,k)-w_1''(C_1,k)|] \le rK\max_{1\le k\le r}\|w_1'-w_1''\|_t.
\]

\item For $H_i$ with $i=3,\ldots,L$, the difference is controlled similarly:
\[
|H_i(W')-H_i(W'')| \le \|L^{(i)}\|_{\infty,1}\max_{1\le k\le r}\E_{C_{i-1}}[|w_{i-1}'(C_{i-1},k)-w_{i-1}''(C_{i-1},k)|] \le rK\max_{1\le k\le r}\|w_{i-1}'-w_{i-1}''\|_t.
\]

\item For $B_k^{(2)}$, the difference is controlled:
\[
|B_k^{(2)}(W')-B_k^{(2)}(W'')| \le K\E_{C_2}[|L^{(2)}_{C_2,k}|(|w_2'(C_2)-w_2''(C_2)| + |w_2''(C_2)||H_2(W')-H_2(W'')|)] \le rK^2(1+B)\|W'-W''\|_t.
\]

\item For $B_k^{(i)}$ with $i=3,\ldots,L-1$, the difference is controlled similarly, giving a factor $(1+B)$ and accumulating factors of $rK$ from each layer.

\item The bad event contributes an exponentially small remainder $e^{-K_1 B^2/2}$.

\item Integrating in time and using Minkowski's inequality yields the result. The factor $(1+rK)^{L-2}$ comes from the accumulation of $rK$ factors from each of the $L-2$ intermediate layers.
\end{enumerate}
\end{proof}

\subsubsection{Complete Proof of Theorem~\ref{thm:wellposed}}

\begin{proof}[Proof of Theorem~\ref{thm:wellposed}]
We perform a Picard-type iteration argument. Consider an arbitrary finite $T\ge 0$ and $W',W''\in\mathcal{W}_T^0$. From Lemma~\ref{lem:difference-F}:
\begin{align*}
\|F(W')-F(W'')\|_t &\le (KK_0(T))^4\left((1+B)\int_0^t \|W'-W''\|_s ds + T\sqrt{2}e^{-K_1 B^2/2}\right)\\
&\equiv k_1(1+B)\int_0^t \|W'-W''\|_s ds + k_2 e^{-k_3 B^2},
\end{align*}
for any $B>0$, where $k_1=(KK_0(T))^4$, $k_2=(KK_0(T))^4 T\sqrt{2}$, and $k_3=K_1/2$.

By Lemma~\ref{lem:invariance-F}, $F$ maps $\mathcal{W}_T^0$ to $\mathcal{W}_T^0$. We can iterate this inequality to obtain:
\begin{align*}
\|F^{(m)}(W')-F^{(m)}(W'')\|_T &\le k_1(1+B)\int_0^T \|F^{(m-1)}(W')-F^{(m-1)}(W'')\|_{T_2} dT_2 + k_2 e^{-k_3 B^2}\\
&\le \frac{1}{m!} T^m k_1^m (1+B)^m \|W'-W''\|_T + k_2 e^{T k_1(1+B) - k_3 B^2}\\
&\le \frac{1}{m!} T^m k_1^m (1+\sqrt{m})^m \|W'-W''\|_T + k_2 e^{T k_1(1+\sqrt{m}) - k_3 m},
\end{align*}
where we choose $B=\sqrt{m}$ in the last display. Note that since $\interleave W\interleave_0<\infty$, $K_0(T)$ and hence $k_1, k_2$ are finite for finite $T$.

By substituting $W''=F(W')$, we obtain:
\[
\sum_{m=1}^\infty \|F^{(m+1)}(W')-F^{(m)}(W')\|_T = \sum_{m=1}^\infty \|F^{(m)}(W'')-F^{(m)}(W')\|_T < \infty.
\]
Hence as $m\to\infty$, $F^{(m)}(W')$ converges in $\|\cdot\|_T$ to a limit $W\in\mathfrak{W}_T$, which is a fixed point of $F$. By Lemma~\ref{lem:bounds-a-priori}, $W$ belongs to $\mathcal{W}_T^0$.

The uniqueness of the fixed point comes from the above estimate, since if $W'$ and $W''$ are fixed points of $F$, then they are both in $\mathcal{W}_T^0$, and:
\[
\|W'-W''\|_T = \|F^{(m)}(W')-F^{(m)}(W'')\|_T \le \frac{1}{m!} T^m k_1^m (1+\sqrt{m})^m \|W'-W''\|_T + k_2 e^{T k_1(1+\sqrt{m}) - k_3 m},
\]
and one can take $m$ arbitrarily large. This proves that the solution exists and is unique on $t\in[0,T]$. Since $T$ is arbitrary, we have existence and uniqueness of the solution to the mean-field ODEs on the time interval $[0,\infty)$.
\end{proof}

\subsection{Proof of ReLU Dense Span Property}
\label{app:relu-dense-span-proof}

This section provides a rigorous proof that frozen random features with full support give dense span in $L^2(\mathcal{P}_X)$, adapting the approach from the full-width case.

\subsubsection{Key Insight: Simpler Than Full-Width Case}

\textbf{Advantage of frozen features:} Unlike the full-width case where we need to prove that full support is \emph{maintained} during training (using algebraic topology), in our case:
\begin{itemize}
\item $f_1(C_1)$ are \textbf{frozen} (not trained), so $\text{supp}(f_1(C_1)) = \mathbb{R}^d$ trivially
\item We only need to prove that full support \textbf{implies} dense span
\item This is a standard universal approximation result, not a dynamical property
\end{itemize}

\subsubsection{Proof Strategy}

The proof proceeds in three steps:

\begin{enumerate}
\item \textbf{Standard universal approximation result:} Use Leshno et al. (1993) to show that $\{\varphi_1(\langle u, \cdot\rangle) : u \in \mathbb{R}^d\}$ has dense span in $L^2(\mathcal{P}_X)$.

\item \textbf{Full support gives dense parameterization:} Since $\text{supp}(f_1(C_1)) = \mathbb{R}^d$, the set $\{f_1(c_1) : c_1 \in \Omega_1\}$ is dense in $\mathbb{R}^d$.

\item \textbf{Continuity argument:} Use continuity of $u \mapsto \varphi_1(\langle u, \cdot\rangle)$ in $L^2(\mathcal{P}_X)$ to show that dense parameterization implies dense span.
\end{enumerate}

\subsubsection{Step 1: Standard Universal Approximation}

\begin{lemma}[Leshno et al. 1993, adapted]
\label{lem:leshno-relu}
Let $\varphi_1$ be ReLU (a non-polynomial activation). Then the function class:
\[
\mathcal{G} = \{\varphi_1(\langle u, \cdot\rangle) : u \in \mathbb{R}^d\}
\]
has dense span in $L^2(\mathcal{P}_X)$.

More precisely, for any $g \in L^2(\mathcal{P}_X)$ and $\epsilon > 0$, there exist $u_1, \ldots, u_n \in \mathbb{R}^d$ and coefficients $\alpha_1, \ldots, \alpha_n \in \mathbb{R}$ such that:
\[
\left\|g - \sum_{i=1}^n \alpha_i \varphi_1(\langle u_i, \cdot\rangle)\right\|_{L^2(\mathcal{P}_X)} < \epsilon.
\]
\end{lemma}

\begin{proof}[Proof sketch]
This follows from the universal approximation theorem of Leshno et al. (1993), which states that single-hidden-layer networks with non-polynomial activations are dense in $C(\mathbb{R}^d)$ (and hence in $L^2$ for compactly supported measures).

For ReLU specifically:
\begin{itemize}
\item ReLU is non-polynomial (it's piecewise linear but not a polynomial)
\item The result extends to $L^2(\mathcal{P}_X)$ when $\mathcal{P}_X$ has compact support (which we assume: $|X| \le K$ with probability 1)
\item The linear combinations $\sum_{i=1}^n \alpha_i \varphi_1(\langle u_i, \cdot\rangle)$ can approximate any function in $L^2(\mathcal{P}_X)$
\end{itemize}

For a complete proof, see Leshno et al. (1993, Theorem 1) or Cybenko (1989) adapted to ReLU.
\end{proof}

\subsubsection{Step 2: Full Support Implies Dense Parameterization}

\begin{lemma}[Full support gives dense parameterization]
\label{lem:dense-parameterization}
If $\text{supp}(f_1(C_1)) = \mathbb{R}^d$, then the set $\{f_1(c_1) : c_1 \in \Omega_1\}$ is dense in $\mathbb{R}^d$.

In other words, for any $u \in \mathbb{R}^d$ and any $\delta > 0$, there exists $c_1 \in \Omega_1$ such that $|f_1(c_1) - u| < \delta$.
\end{lemma}

\begin{proof}
This is immediate from the definition of support. If $\text{supp}(f_1(C_1)) = \mathbb{R}^d$, then for any open ball $B_\delta(u) \subset \mathbb{R}^d$, we have:
\[
\mathbb{P}(f_1(C_1) \in B_\delta(u)) > 0.
\]
This means there exists at least one $c_1 \in \Omega_1$ such that $f_1(c_1) \in B_\delta(u)$, i.e., $|f_1(c_1) - u| < \delta$.

Since this holds for any $u \in \mathbb{R}^d$ and any $\delta > 0$, the set $\{f_1(c_1) : c_1 \in \Omega_1\}$ is dense in $\mathbb{R}^d$.
\end{proof}

\subsubsection{Step 3: Continuity Argument}

\begin{lemma}[Continuity of ReLU features]
\label{lem:continuity-relu}
The mapping $u \mapsto \varphi_1(\langle u, \cdot\rangle)$ is continuous from $\mathbb{R}^d$ to $L^2(\mathcal{P}_X)$.

More precisely, for any $u, u' \in \mathbb{R}^d$:
\[
\|\varphi_1(\langle u, \cdot\rangle) - \varphi_1(\langle u', \cdot\rangle)\|_{L^2(\mathcal{P}_X)} \le K |u - u'|,
\]
where $K$ depends on the bound on $|X|$ and the Lipschitz constant of $\varphi_1$.
\end{lemma}

\begin{proof}
Since $\varphi_1$ is ReLU (which is 1-Lipschitz) and $|X| \le K$ with probability 1, we have:
\begin{align*}
\|\varphi_1(\langle u, \cdot\rangle) - \varphi_1(\langle u', \cdot\rangle)\|_{L^2(\mathcal{P}_X)}^2 
&= \mathbb{E}_X[|\varphi_1(\langle u, X\rangle) - \varphi_1(\langle u', X\rangle)|^2]\\
&\le \mathbb{E}_X[|\langle u - u', X\rangle|^2] \quad \text{(ReLU is 1-Lipschitz)}\\
&\le |u - u'|^2 \mathbb{E}_X[|X|^2]\\
&\le K^2 |u - u'|^2.
\end{align*}
Taking square roots gives the desired bound.
\end{proof}

\subsubsection{Completing the Proof of Theorem~\ref{thm:relu-dense-span}}

\begin{proof}[Proof of Theorem~\ref{thm:relu-dense-span}]
By Lemma~\ref{lem:leshno-relu}, for any $g \in L^2(\mathcal{P}_X)$ and $\epsilon > 0$, there exist $u_1, \ldots, u_n \in \mathbb{R}^d$ and coefficients $\alpha_1, \ldots, \alpha_n$ such that:
\[
\left\|g - \sum_{i=1}^n \alpha_i \varphi_1(\langle u_i, \cdot\rangle)\right\|_{L^2(\mathcal{P}_X)} < \frac{\epsilon}{2}.
\]

By Lemma~\ref{lem:dense-parameterization}, since $\text{supp}(f_1(C_1)) = \mathbb{R}^d$, for each $u_i$ and $\delta > 0$ (to be chosen), there exists $c_i \in \Omega_1$ such that $|f_1(c_i) - u_i| < \delta$.

By Lemma~\ref{lem:continuity-relu}, we have:
\[
\|\varphi_1(\langle f_1(c_i), \cdot\rangle) - \varphi_1(\langle u_i, \cdot\rangle)\|_{L^2(\mathcal{P}_X)} \le K |f_1(c_i) - u_i| < K\delta.
\]

Therefore:
\begin{align*}
\left\|g - \sum_{i=1}^n \alpha_i \varphi_1(\langle f_1(c_i), \cdot\rangle)\right\|_{L^2(\mathcal{P}_X)}
&\le \left\|g - \sum_{i=1}^n \alpha_i \varphi_1(\langle u_i, \cdot\rangle)\right\|_{L^2(\mathcal{P}_X)}\\
&\quad + \left\|\sum_{i=1}^n \alpha_i (\varphi_1(\langle u_i, \cdot\rangle) - \varphi_1(\langle f_1(c_i), \cdot\rangle))\right\|_{L^2(\mathcal{P}_X)}\\
&< \frac{\epsilon}{2} + \sum_{i=1}^n |\alpha_i| K\delta.
\end{align*}

Choosing $\delta = \frac{\epsilon}{2K \sum_{i=1}^n |\alpha_i|}$ (which is finite since we can choose the $\alpha_i$ to be bounded), we get:
\[
\left\|g - \sum_{i=1}^n \alpha_i \varphi_1(\langle f_1(c_i), \cdot\rangle)\right\|_{L^2(\mathcal{P}_X)} < \epsilon.
\]

This proves that $\mathcal{F}$ has dense span in $L^2(\mathcal{P}_X)$.
\end{proof}

\subsection{Proof of Channel Specialization Theorem}
\label{app:r2-logratio-proof}

\begin{proof}[Proof of Theorem~\ref{thm:r2-logratio-criterion-main}]
By the exact identity for the log-ratio derivative at $x=x_0$,
\[
\partial_t R_{12}(t,x_0)=\frac{\mathrm{sign}(f_1(t,x_0))\,\partial_t f_1(t,x_0)}{|f_1(t,x_0)|+\varepsilon}
-\frac{\mathrm{sign}(f_2(t,x_0))\,\partial_t f_2(t,x_0)}{|f_2(t,x_0)|+\varepsilon}.
\]
Inserting the local dynamics $\partial_t f_k(t,x_0)=-\xi_1(t)\,K_{\mu_0}(x_0,x_0)\,d_0(t)\,B_{k,1}(t)$ and using $-d_0(t)\ge 0$:
\[
\partial_t R_{12}(t,x_0)=\xi_1(t)\,K_{\mu_0}(x_0,x_0)\,(-d_0(t))\Big(
\frac{\mathrm{sign}(f_1(t,x_0))B_{1,1}(t)}{|f_1(t,x_0)|+\varepsilon}-\frac{\mathrm{sign}(f_2(t,x_0))B_{2,1}(t)}{|f_2(t,x_0)|+\varepsilon}\Big).
\]
By the sign assumption on $B_{1,1}$, we have $\mathrm{sign}(f_1(t,x_0))B_{1,1}(t)=|B_{1,1}(t)|$. Moreover, $-\mathrm{sign}(f_2(t,x_0))B_{2,1}(t)\ge -|B_{2,1}(t)|$. Using the stability assumption \eqref{eq:r2-stability-assumption-main} yields
\[
\frac{|B_{1,1}(t)|}{|f_1(t,x_0)|+\varepsilon}-\frac{\mathrm{sign}(f_2(t,x_0))B_{2,1}(t)}{|f_2(t,x_0)|+\varepsilon}
\ge \frac{|B_{1,1}(t)|}{|f_1(t,x_0)|+\varepsilon}-\frac{|B_{2,1}(t)|}{|f_2(t,x_0)|+\varepsilon}
\ge (1-\rho)\frac{|B_{1,1}(t)|}{|f_1(t,x_0)|+\varepsilon},
\]
which proves the claimed lower bound for $\partial_t R_{12}$.
\end{proof}

\subsection{Proof of Quantitative Approximation}
\label{app:quantitative-proof}

This section provides a rigorous quantitative bound on the approximation error between finite-width low-rank networks and their mean-field limit. The key result shows that the approximation error scales as $O(1/\sqrt{n_{\min}} + \sqrt{\epsilon})$ with explicit constants, where $n_{\min}$ is the minimum width across layers and $\epsilon$ is the learning rate step size.

\subsubsection{Particle ODEs for Low-Rank Networks}

We construct auxiliary trajectories, which we call the \textit{particle ODEs} for the low-rank case. These are continuous-time trajectories of finitely many neurons, averaged over the data distribution, adapted to the low-rank structure:
\begin{align*}
\frac{\partial}{\partial t}\tilde{w}_2(t, j_2) &= -\xi_2(t)\mathbb{E}_Z\left[d_L(Z; \tilde{W}(t))\,\varphi_2(\tilde{H}_2(t, j_2; X, \tilde{W}(t)))\right],\\
\frac{\partial}{\partial t}\tilde{w}_1(t, j_1, k) &= -\xi_1(t)\mathbb{E}_Z\left[d_L(Z; \tilde{W}(t))\,\varphi_1(\langle f_1(C_1(j_1)), X\rangle)\,\tilde{B}_k(t; X, \tilde{W}(t))\right],
\end{align*}
where $j_1 = 1, \ldots, n_1$, $j_2 = 1, \ldots, n_2$, $k = 1, \ldots, r$, $\tilde{W}(t) = (\tilde{w}_1(t, \cdot, \cdot), \tilde{w}_2(t, \cdot))$, and $t \in \mathbb{R}_{\ge 0}$. 

The second-layer output for the particle ODEs is:
\[
\tilde{H}_2(t, j_2; X, \tilde{W}) = \sum_{k=1}^r L_{C_2(j_2),k}\,\frac{1}{n_1}\sum_{j_1=1}^{n_1} \tilde{w}_1(t, j_1, k)\,\varphi_1(\langle f_1(C_1(j_1)), X\rangle),
\]
and the backpropagated signal is:
\[
\tilde{B}_k(t; X, \tilde{W}) = \frac{1}{n_2}\sum_{j_2=1}^{n_2} L_{C_2(j_2),k}\,\varphi_2'(\tilde{H}_2(t, j_2; X, \tilde{W}))\,\tilde{w}_2(t, j_2).
\]

We specify the initialization $\tilde{W}(0)$: $\tilde{w}_1(0, j_1, k) = w_1^0(C_1(j_1), k)$ and $\tilde{w}_2(0, j_2) = w_2^0(C_2(j_2))$. That is, it shares the same initialization with the neural network $\mathbf{W}(0)$, and hence is coupled with the neural network and the MF ODEs.

\subsubsection{Key Lemmas}

\begin{lemma}[Particle coupling bound]
\label{lem:particle-lowrank}
Under the same setting as Theorem~\ref{thm:quantitative-lowrank}, for any $\delta > 0$, with probability at least $1-\delta$,
\[
\mathscr{D}_T(W, \tilde{W}) \le \frac{1}{\sqrt{n_{\min}}}\log^{1/2}\left(\frac{3(T+1)n_{\max}^2r^{L-2}}{\delta} + e\right)\mathbf{e^{K_T\textcolor{red}{(1+rK)^{L-2}}}},
\]
in which $n_{\min} = \min\{n_1, \ldots, n_{L-1}\}$, $n_{\max} = \max\{n_1, \ldots, n_{L-1}\}$, and $K_T = K(1+T^K)$.
\end{lemma}

\begin{proof}
The proof follows the same structure as the full-width case, with the key adaptation being the low-rank structure. We decompose the differences $q_i$ and $q_{B,k}$ as described in the main text, apply concentration bounds with union bounds over $r$ channels, and use Gronwall's lemma. The exponential constant accumulates as $(1+rK)^{L-2}$ because each intermediate layer contributes a factor $(1+rK)$, and these factors multiply across layers. The detailed steps are provided in the main text (Section~\ref{rem:generalization-L}).
\end{proof}

\begin{lemma}[Gradient descent discretization error]
\label{lem:gradient-lowrank}
Under the same setting as Theorem~\ref{thm:quantitative-lowrank}, for any $\delta > 0$ and $\epsilon \le 1$, with probability at least $1-\delta$,
\[
\mathscr{D}_T(\tilde{W}, \mathbf{W}) \le \sqrt{\epsilon\log\left(\frac{\mathbf{2n_1n_2r^{L-2}}}{\delta} + e\right)}\mathbf{e^{K_T\textcolor{red}{(1+rK)^{L-2}}}},
\]
in which $K_T = K(1+T^K)$.
\end{lemma}

\begin{proof}
The proof follows the same structure as the full-width case, with careful adaptations for the low-rank structure. The key difference is that we need to account for the $r$ channels in $w_1$ and the mixing matrix $L$ in all bounds. The martingale concentration bounds include factors of $r$ from the union bound over channels, and the exponential constant accumulates as $(1+rK)^{L-2}$ for $L$ layers. The detailed proof is provided in the main text.
\end{proof}

\begin{proof}[Proof of Theorem~\ref{thm:quantitative-lowrank}]
Using the triangle inequality:
\[
\mathscr{D}_T(W, \mathbf{W}) \le \mathscr{D}_T(W, \tilde{W}) + \mathscr{D}_T(\tilde{W}, \mathbf{W}),
\]
the result follows immediately from Lemmas~\ref{lem:particle-lowrank} and \ref{lem:gradient-lowrank}, noting that the $\log(2n_1n_2r^{L-2}/\delta)$ term can be absorbed into the $\log(3(T+1)n_{\max}^2r^{L-2}/\delta + e)$ term up to constants.
\end{proof}

\subsection{Proof of Global Convergence}
\label{app:global-convergence-proof}

We now provide the detailed proof, adapting the argument from the full-width three-layer case to our low-rank setting. The key simplification is that we do not train $\rho^1$ (the first layer uses frozen random features), so we avoid the algebraic topology arguments needed to show full support of $\text{Law}(w_1^*(t,C_1))$. $H_i$ bi-Lipschitz in $W$ is sufficient and equivalent to $\varphi_i$ Lipschitz; it is the only thing that matters. After defining norms and updating $K_0(t)$ with $(1+rK)^{1/2}$, the remaining is simple; all steps are in one proof below.

\begin{proof}[Proof of Theorem~\ref{thm:global-convergence-lowrank}]

\textbf{Step 1: Convergence of the backpropagated signal}

By the convergence assumption (Assumption~\ref{assump:convergence-lowrank}), we have that for any $\epsilon>0$, there exists $T(\epsilon)$ such that for all $t\ge T(\epsilon)$ and almost surely:
\[
\epsilon \ge \left|\E_{Z}\!\left[d_L\big(Z;W^*(t)\big)\,\varphi_1\!\left(\ip{f_1(C_1)}{X}\right)\,B_k^{(2)*}(t;X)\right]\right| = \left|\left\langle \E_{Z}\!\left[d_L\big(Z;W^*(t)\big)\,B_k^{(2)*}(t;X)\middle|X=x\right],\varphi_1(\ip{f_1(C_1)}{x})\right\rangle_{L^2(\mathcal{P}_X)}\right|,
\]
where $B_k^{(2)*}(t;X)$ is the backpropagated signal at the solution $W^*(t)$.

\textbf{Step 2: Universal approximation property}

Since $f_1(C_1)$ are frozen random features with full support in $\R^d$ (by Assumption~\ref{assump:diversity-lowrank}), and we have $r$ channels, the set $\{\varphi_1(\ip{f_1(c_1)}{\cdot}): c_1\in\Omega_1, k\in\{1,\dots,r\}\}$ has dense span in $L^2(\mathcal{P}_X)$ by Theorem~\ref{thm:relu-dense-span}. We thus obtain that for $u$ in a dense subset of $\R^d$:
\[
\text{ess-sup}\left|\left\langle \E_{Z}\!\left[d_L\big(Z;W^*(t)\big)\,B_k^{(2)*}(t;X)\middle|X=x\right],\varphi_1(\ip{u}{x})\right\rangle_{L^2(\mathcal{P}_X)}\right| \le \epsilon.
\]

By continuity of $u\mapsto\varphi_1(\ip{u}{\cdot})$ in $L^2(\mathcal{P}_X)$, we extend the above to all $u\in\R^d$.

\textbf{Step 3: Coupling argument with low-rank structure}

Recall the couplings $\pi_t$ in Assumption~\ref{assump:convergence-lowrank}. Since $\varphi_1$ is bounded, we have:
\begin{align*}
&\E_{(C_1,C_2,C_1',C_2')\sim\pi_t}\left[\left|\left\langle \E_{Z}\!\left[d_L\big(Z;W^*(t)\big)\,B_k^{(2)*}(t;X)-d_L\big(Z;\bar{W}\big)\,\bar{B}_k^{(2)}(X)\middle|X=x\right],\varphi_1(\ip{u}{x})\right\rangle_{L^2(\mathcal{P}_X)}\right|\right]\\
&\quad \le K\,\E_{\pi_t}\left[\left|d_L\big(Z;W^*(t)\big)\,B_k^{(2)*}(t;X)-d_L\big(Z;\bar{W}\big)\,\bar{B}_k^{(2)}(X)\right|\right],
\end{align*}
where $\bar{W}=(\bar{w}_1,\ldots,\bar{w}_L)$ and $\bar{B}_k^{(2)}(X)$ is the backpropagated signal at the limit point.

\textbf{Step 4: Key bound using bi-Lipschitz property}

Using the regularity assumption and the low-rank structure, similar to the calculation in the full-width case, we bound:
\begin{align*}
&\E_{\pi_t}\left[\left|d_L\big(Z;W^*(t)\big)\,B_k^{(2)*}(t;X)-d_L\big(Z;\bar{W}\big)\,\bar{B}_k^{(2)}(X)\right|\right]\\
&\quad \le K\,\E_{\pi_t}\Big[(1+|\bar{w}_L(C_L)|)\prod_{j=2}^{L-1}(1+|\bar{w}_j(C_j)|)\Big(|w_L^*(t,C_L')-\bar{w}_L(C_L)|\\
&\qquad + \sum_{j=2}^{L-1}\max_{1\le k\le r}|w_j^*(t,C_j',k)-\bar{w}_j(C_j,k)|\Big)\Big],
\end{align*}
where the last step uses:
\begin{itemize}
\item The Lipschitz property of $\partial_2\mathcal{L}$ and $\varphi_i'$
\item \textcolor{red}{The bi-Lipschitz bounds for $H_i$ differences with low-rank structure}
\item \textcolor{red}{The structure $B_k^{(i)} = \E_{C_i}[L^{(i)}_{C_i,k}\,\varphi_i'(H_i)\,w_i]$, which gives bounds involving factors of $rK$}
\end{itemize}

By Assumption~\ref{assump:convergence-lowrank}, the right-hand side converges to $0$ as $t\to\infty$.

\textbf{Step 5: Orthogonality and optimality}

We thus obtain that for all $u\in\R^d$:
\[
\E_{C_L}\left[\left|\left\langle \E_{Z}\!\left[d_L\big(Z;\bar{W}\big)\,\bar{B}_k^{(2)}(X)\middle|X=x\right],\varphi_1(\ip{u}{x})\right\rangle_{L^2(\mathcal{P}_X)}\right|\right] = 0,
\]
which yields that for all $u\in\R^d$ and almost surely:
\[
\left|\left\langle \E_{Z}\!\left[d_L\big(Z;\bar{W}\big)\,\bar{B}_k^{(2)}(X)\middle|X=x\right],\varphi_1(\ip{u}{x})\right\rangle_{L^2(\mathcal{P}_X)}\right| = 0.
\]

Since $\{\varphi_1(\ip{u}{\cdot}): u\in\R^d\}$ has dense span in $L^2(\mathcal{P}_X)$ (by Theorem~\ref{thm:relu-dense-span}), we have:
\[
\E_{Z}\!\left[d_L\big(Z;\bar{W}\big)\,\bar{B}_k^{(2)}(X)\middle|X=x\right] = 0
\]
for $\mathcal{P}_X$-almost every $x$ and almost every $c_2$.

Expanding $\bar{B}_k^{(2)}(X)$ and using the structure of the backpropagated signals through all layers, we obtain that the gradient is zero at the limit point.

\textbf{Step 6: Non-degeneracy}

We ensure that the limit point is non-degenerate (i.e., $\Pp(\bar{w}_L(C_L) \ne 0) > 0$). This is crucial for the optimality argument. There are two cases:

\begin{itemize}
\item \textbf{Case A:} If the output layer is frozen ($\xi_L(\cdot) = 0$) and the initialization has non-zero mass, then non-degeneracy is automatic since the weights don't change.

\item \textbf{Case B:} If the output layer is trained ($\xi_L(\cdot) \ne 0$), we use Assumption~\ref{assump:nondegeneracy-lowrank}: $\mathscr{L}(w_1^0,\ldots,w_L^0) < \E_Z[\mathcal{L}(Y,\varphi_L(0))]$.

By the gradient flow property (loss decreases along trajectories):
\[
\mathscr{L}(w_1^*(t),\ldots,w_L^*(t)) \le \mathscr{L}(w_1^*(t'),\ldots,w_L^*(t')) \quad \text{for } t\ge t'.
\]
In particular, setting $t'=0$ and taking $t\to\infty$:
\[
\mathscr{L}(\bar{w}_1,\ldots,\bar{w}_L) \le \mathscr{L}(w_1^0,\ldots,w_L^0) < \E_Z[\mathcal{L}(Y,\varphi_L(0))].
\]

Now, if the limit point were degenerate (i.e., $\Pp(\bar{w}_L(C_L) = 0) = 1$), then the network output would be:
\[
\hat{y}(X; \bar{W}) = \E_{C_L}[\bar{w}_L(C_L)\,\varphi_L(H_L(C_L; X, \bar{W}))] = 0,
\]
since $\bar{w}_L = 0$ almost surely. This means $\mathscr{L}(\bar{w}_1,\ldots,\bar{w}_L) = \E_Z[\mathcal{L}(Y, \varphi_L(0))]$, which contradicts the inequality above. Therefore, $\Pp(\bar{w}_L(C_L) \ne 0) > 0$, establishing non-degeneracy.
\end{itemize}

Since $\varphi_i'$ is strictly non-zero and the mixing kernels $L^{(i)}$ have non-zero entries, we have:
\[
\E_{Z}\!\left[\partial_2\mathcal{L}\big(Y,\hat y(X;\bar{W})\big)\middle|X=x\right] = 0
\]
for $\mathcal{P}_X$-almost every $x$.

\textbf{Step 7: Global optimality}

Since $\mathcal{L}$ is convex in its second variable, for any measurable function $\tilde{y}(x)$:
\[
\mathcal{L}\big(y,\tilde{y}(x)\big)-\mathcal{L}\big(y,\hat y(x;\bar{W})\big) \ge \partial_2\mathcal{L}\big(y,\hat y(x;\bar{W})\big)\big(\tilde{y}(x)-\hat y(x;\bar{W})\big).
\]

Taking expectation:
\[
\E_Z\big[\mathcal{L}\big(Y,\tilde{y}(X)\big)\big] \ge \mathscr{L}(\bar{w}_1,\ldots,\bar{w}_L),
\]
i.e., $(\bar{w}_1,\ldots,\bar{w}_L)$ is a global minimizer.

\textbf{Step 8: Connection to the limit}

Finally, to connect $\mathscr{L}(W^*(t))$ with $\mathscr{L}(\bar{w}_1,\ldots,\bar{w}_L)$ in the limit $t\to\infty$, we have:
\begin{align*}
\big|\mathscr{L}(W^*(t))-\mathscr{L}(\bar{w}_1,\ldots,\bar{w}_L)\big| &= \left|\E_Z\big[\mathcal{L}\big(Y,\hat y^*(t,X)\big)-\mathcal{L}\big(Y,\hat y(X;\bar{W})\big)\big]\right|\\
&\le K\,\E_Z\big[\big|\hat y^*(t,X)-\hat y(X;\bar{W})\big|\big]\\
&\le K\,\E_{\pi_t}\Big[(1+|\bar{w}_L(C_L)|)\prod_{j=2}^{L-1}(1+|\bar{w}_j(C_j)|)\Big(|w_L^*(t,C_L')-\bar{w}_L(C_L)|\\
&\quad + \sum_{j=2}^{L-1}\max_{1\le k\le r}|w_j^*(t,C_j',k)-\bar{w}_j(C_j,k)|\Big)\Big],
\end{align*}
where the last step uses the Lipschitz property of $\varphi_L$ and \textcolor{red}{the bi-Lipschitz bounds for $H_i$ with low-rank structure}.

By Assumption~\ref{assump:convergence-lowrank}, the right-hand side converges to $0$ as $t\to\infty$. This completes the proof.
\end{proof}

\end{document}
